\chapter{Implementation}
\label{ch:implementation}

This chapter describes the realised system: the technology stack and the
reasoning behind each choice, the service interface, the two web interfaces, the
research database, the deployment modes, and the reproducibility measures.

\section{Technology Stack}
\label{sec:stack}

Table~\ref{tab:stack} records each major technology, where it is used, why it
was selected, and its principal drawback. The drawbacks are stated explicitly
because several of them constrain claims made later in this report.

\begin{table}[H]
\centering
\caption{Technology stack with rationale and trade-offs.}
\label{tab:stack}
\scriptsize
\begin{tabularx}{\textwidth}{@{}L{0.92in} L{0.85in} X X@{}}
\toprule
\textbf{Technology} & \textbf{Where used} & \textbf{Why chosen} &
\textbf{Drawback} \\
\midrule
scikit-learn \cite{pedregosa2011} & Cascade training and inference & Mature,
interpretable, strong on tabular clinical data; fast CPU inference &
Lower ceiling than deep tabular methods on very large data \\
imbalanced-learn & Training-time resampling & Standard implementations of
over/under-sampling &
Risk of overfitting synthetic minority patterns if not cross-validated \\
joblib & Model serialisation & Plain-\texttt{dict} package with no custom
classes; self-describing; service trusts the package's own feature list &
Python-specific format \\
FastAPI \cite{fastapi} + Uvicorn & Backend service & Async, automatic OpenAPI
documentation, native Pydantic validation &
Blocking ML/LLM calls must be managed carefully in an async context \\
Pydantic & Request validation & \texttt{extra="forbid"} plus per-field ranges
reject malformed input before the model sees it &
Schema must track the model's feature list \\
BGE-M3 via Ollama & RAG embedding (1024-d) & Strong open multilingual embedder;
serving via Ollama avoids a \texttt{torch.load} version restriction &
Requires the Ollama service to be running \\
FAISS & Vector search & Fast, standard, integrates cleanly &
Overkill at 162 passages; value is headroom for a larger corpus \\
Qwen3-4B-Instruct-2507 & Text generation & Strong instruction-following at a
size that is fine-tunable and self-hostable; no patient data leaves the machine &
Quality ceiling below frontier models; needs a GPU for comfortable latency \\
PEFT & LoRA fine-tuning (bf16 base) & Trains a 4B model on a single GPU;
adapter is small and swappable &
Adapter quality bounded by the 540-example corpus \\
Safety validator (pure Python) & Post-generation gate & Deterministic,
auditable, zero inference cost &
Cannot cover unseen unsafe phrasings \\
PostgreSQL + SQLAlchemy + Alembic & Research persistence & JSONB for nested
records, typed columns for queries, versioned migrations &
An extra service; absence is silent unless health is checked \\
React + Vite + Tailwind & Web interfaces & Fast builds, small bundle,
no heavyweight state library needed &
No server-side rendering or routing framework \\
\bottomrule
\end{tabularx}
\end{table}

\subsection{A note on the embedding runtime}

During development the retrieval index was found to have been built with a
lexical hashing fallback rather than the intended BGE-M3 embedder. The cause was
a library-level security restriction: the installed PyTorch version predates the
floor required for a safe \texttt{torch.load} path, so the
\texttt{sentence-transformers} loader raised, and an automatic embedder-selection
routine caught the exception and silently degraded to the fallback, emitting only
a printed warning.

Two remedies were available: upgrade PyTorch, or serve the same embedding model
through a different runtime. The second was chosen, because the language-model
serving stack was already verified working against the installed version and an
upgrade risked destabilising it for no modelling benefit. BGE-M3 is therefore
served through Ollama, which does not use the restricted loading path. The index
was rebuilt, and the retrieval evaluation reported in Chapter~\ref{ch:results}
was re-run entirely against the corrected index.

The episode is documented here rather than omitted because it is
methodologically relevant: a silent fallback of this kind degrades results
without producing an error, and the corrected and uncorrected runs were both
re-scored so that the difference could be attributed.

\section{Backend Service}
\label{sec:backend}

The backend exposes the endpoints in Table~\ref{tab:api}. The model is loaded
once at startup; the language-model runtime is a lazily initialised singleton so
that a deployment which never requests guidance never pays the model-loading
cost.

\begin{table}[H]
\centering
\caption{HTTP API surface.}
\label{tab:api}
\small
\begin{tabularx}{\textwidth}{@{}l L{1.85in} X@{}}
\toprule
\textbf{Method} & \textbf{Endpoint} & \textbf{Purpose} \\
\midrule
GET & \texttt{/health} & Service and model status, library versions, database state \\
GET & \texttt{/db/health} & Database connectivity and stored record count \\
GET & \texttt{/features} & Feature schema, metadata, numeric and binary lists \\
GET & \texttt{/model-info} & Thresholds and per-stage feature lists \\
POST & \texttt{/predict} & Cascade scoring: category, probabilities,
explanation, reliability \\
POST & \texttt{/reload-model} & Hot-reload the serialised model package \\
GET & \texttt{/generate-guidance/health} & Language-model mode, model
specification, active prompt variant \\
POST & \texttt{/generate-guidance} & Full pipeline: predict, retrieve,
generate, validate, cite \\
\bottomrule
\end{tabularx}
\end{table}

Two design details support auditability. First, the model specification string
reported by the health endpoint encodes adapter status directly---for example
\texttt{qwen:Qwen/Qwen3-4B-Instruct-2507+lora:v1/adapter}---so the question
``is the fine-tuned model actually serving this request?'' is answerable from a
single response, and from any stored run record, without inspecting the process
that produced it. Second, \texttt{/generate-guidance} returns the entire audit
trail alongside the patient-facing text: the retrieval query, the retrieved
passage identifiers with similarity scores, the raw pre-validator text, the
validator status and the list of violations, and the complete prediction object.

\section{Web Interfaces}
\label{sec:frontend}

The frontend is a single-page React application presenting two distinct
interfaces over the same backend.

The \textbf{patient interface} presents a fourteen-field form grouped into
sections, with per-field normal ranges, inline validation and a live
completeness indicator. The result view shows the risk category with a plain
statement that it is a screening result and not a diagnosis, the reliability
level with an input-completeness bar, a row of ``main signals'' derived from the
top attributed factors expressed in everyday language, and the five education
sections with inline citation markers and a reference list.

The \textbf{research interface} exposes the full audit trail: both stage
probabilities and thresholds, the normalised class probabilities, the input
quality panel listing imputed and out-of-range values, the complete factor
attribution table with contribution magnitudes, the retrieval panel with the
query string and per-passage scores and topics, a side-by-side view of raw and
post-validator text, the safety panel showing which validator layers fired, a
live system-health panel, and the frozen evaluation metrics.

Exposing the raw pre-validator text next to the delivered text was a deliberate
decision. It makes the guardrail's effect directly observable during a
demonstration---one can see precisely what the model produced and what was
delivered instead---rather than asking an audience to take the safety claim on
trust.

\section{Research Database}
\label{sec:database}

Every screening is recorded in a PostgreSQL table, whose schema is given in
Table~\ref{tab:dbschema}. Nested structures (raw inputs, full prediction, full
guidance response) are stored as JSONB so that the schema does not require
migration each time a field is added, while the fields most useful for research
queries are additionally materialised as typed, indexed columns.

\begin{table}[H]
\centering
\caption{The \texttt{screenings} table.}
\label{tab:dbschema}
\small
\begin{tabularx}{\textwidth}{@{}L{1.55in} L{1.15in} X@{}}
\toprule
\textbf{Column} & \textbf{Type} & \textbf{Contents} \\
\midrule
\texttt{id} & INTEGER PK & Row identifier \\
\texttt{created\_at} & TIMESTAMPTZ & Time of screening (indexed) \\
\texttt{source} & VARCHAR(32) & \texttt{predict} or \texttt{generate-guidance} (indexed) \\
\texttt{input\_features} & JSONB & Raw submitted values \\
\texttt{n\_missing\_features} & INTEGER & Count of imputed inputs \\
\texttt{risk\_class} & VARCHAR(16) & low / moderate / high (indexed) \\
\texttt{risk\_class\_display} & VARCHAR(64) & Human-readable category label \\
\texttt{confidence} & FLOAT & Maximum class probability \\
\texttt{probabilities} & JSONB & Full probability vector \\
\texttt{prediction} & JSONB & Complete \texttt{/predict} response \\
\texttt{guidance} & JSONB & Complete guidance response, when applicable \\
\texttt{safety\_status} & VARCHAR(16) & pass / sanitized / fallback \\
\texttt{llm\_mode} & VARCHAR(32) & production / local\_demo \\
\texttt{model\_spec} & VARCHAR(128) & Model identifier including adapter status \\
\bottomrule
\end{tabularx}
\end{table}

Persistence is best-effort by design. The database layer initialises lazily and,
if the connection fails, every helper degrades to a no-op while the API
continues to serve screenings normally. A failed write is logged, never
surfaced to the user, and never allowed to fail a request. Schema is managed
explicitly through Alembic migrations rather than automatic table creation, so
that the deployed schema is version-controlled.

\section{Deployment Modes}
\label{sec:modes}

Three interchangeable modes are supported behind one interface, selected by
environment variable and reported by the health endpoint
(Table~\ref{tab:modes}).

\begin{table}[H]
\centering
\caption{Deployment modes.}
\label{tab:modes}
\scriptsize
\begin{tabularx}{\textwidth}{@{}L{0.85in} X X X@{}}
\toprule
& \textbf{local\_demo} & \textbf{ollama} & \textbf{qwen (deployed)} \\
\midrule
Language model & Deterministic mock generator & Qwen3-4B via Ollama REST &
Qwen3-4B via \texttt{transformers} + LoRA adapter \\
Embedder & Hashing TF-IDF (4096-d) & BGE-M3 (1024-d) & BGE-M3 (1024-d) \\
Vector search & Brute force & Brute force / FAISS & FAISS \\
Prompt variant & Full & Full & Compact (adapter active) \\
GPU required & No & No & Recommended \\
Safety validator & Identical & Identical & Identical \\
Purpose & Offline demonstration, CI & Laptop GPU without the ML stack &
The evaluated configuration \\
\bottomrule
\end{tabularx}
\end{table}

The offline mode exists so that the complete pipeline---retrieval, prompt
construction, validation, citation---can be exercised on any machine without a
GPU or network. Its generator is deterministic and is never used for any
reported evaluation figure. The safety validator is identical across all three
modes, which means the guardrail behaviour can be demonstrated even where the
production model cannot run.

\section{Reproducibility}
\label{sec:repro}

Several measures support reproduction of the reported results.

\begin{itemize}
  \item \textbf{Dataset integrity.} The prepared modelling file is
        hash-verified, and the split audit is re-run and asserted at training
        time.
  \item \textbf{Configuration fingerprinting.} Cross-validation seeds, model
        seed, bootstrap replicate count, evaluation seeds and thresholds are
        recorded together with a configuration fingerprint.
  \item \textbf{Serialisation verification.} After the model package is written,
        it is reloaded and its predictions compared against the in-memory model;
        agreement was exact to within $2.2\times10^{-16}$, and the package
        checksum is recorded.
  \item \textbf{Frozen evaluation benchmark.} The 300-profile language-model
        benchmark is content-fingerprinted, and the fingerprint was re-verified
        before and after every generation run to confirm that all arms were
        evaluated against an identical benchmark.
  \item \textbf{Append-only run records.} Every generated response is stored
        with the hash of its own raw text, allowing a later re-scoring to prove
        that only scoring changed and generation did not.
  \item \textbf{Deterministic decoding.} All reported language-model results use
        a fixed seed with deterministic decoding.
\end{itemize}

The deployed model package is 302~KB, which together with the 132~MB adapter and
the 162-passage index constitutes the entire deployable artefact set beyond the
base model weights.
